% ============ Chapter 7. Regular rings (原书 pp.153-173) ============
\chapter{Regular rings}\label{ch:7}

Regular local rings have already been mentioned several times, and in
this chapter we are going to carry out a study of them using homological
algebra. Serre's Theorem~\ref{thm:19.2}, characterising regular local rings as Noetherian
local rings of finite global dimension, is the really essential result, and
from this one can deduce at once, for example, that a localisation of a
regular local ring is again regular (Theorem~\ref{thm:19.3}); this is a result which
ideal theory on its own was only able to prove with difficulty in special
cases. \S\ref{sec:20} on UFDs is centred around the theorem that a regular local
ring is a UFD, another important achievement of homological methods;
we only cover the basic topics. This section was written referring to the early
parts of Professor M.\ Narita's lectures at Tokyo Metropolitan University.
In \S\ref{sec:21} we give a simple discussion of the most elementary results on
complete intersection rings. This is an area where the homology theory of
rings plays an essential role, but we are only able to mention this in
passing.

% ============ §19 Regular rings (原书 pp.153-161) ============
\section{Regular rings}\label{sec:19}\index{regular local rings, section 19}\index{regular rings, section 19}\index{minimal resolution, section 19}\index{global dimension, section 19}

\textbf{Minimal free resolutions.} Let $(A, m, k)$ be a local ring, $M$ and $N$ finite
$A$-modules. An $A$-linear map $\varphi\colon M \to N$ induces a $k$-linear map
$\bar{\varphi}\colon M \otimes k \to N \otimes k$, which we denote $\bar{\varphi}$; then one sees easily that
\[
  \varphi \text{ is an isomorphism } \iff \bar{\varphi} \text{ is surjective and } \Ker \varphi = mM.
\]
In particular for free modules $M$ and $N$, if $\bar{\varphi}$ is an isomorphism then $\rank
M = \rank N$, and writing $\varphi$ as a matrix we have $\det \varphi \notin m$, so that
\[
  \varphi \text{ is an isomorphism } \iff \bar{\varphi} \text{ is an isomorphism}.
\]
Let $M$ be a finite $A$-module. An exact sequence
\[
  (*) \quad \cdots \xrightarrow{d_{i+1}} L_i \xrightarrow{d_i} L_{i-1} \to \cdots \to L_1 \xrightarrow{d_1} L_0 \xrightarrow{\varepsilon} M \to 0
\]
(or the complex $L_\bullet$) is called a \emph{minimal (free) resolution}\index{minimal resolution} of $M$ if it satisfies the
three conditions (1) each $L_i$ is a finite free $A$-module, (2) $\bar{d}_i = 0$, or in other
words $d_iL_i \subset mL_{i-1}$ for all $i$, and (3) $\bar{\varepsilon}\colon L_0 \otimes k \to M \otimes k$ is an isomorph-
ism. Breaking up $(*)$ into short exact sequences $0 \to K_1 \to L_0 \to M \to 0$,
$0 \to K_2 \to L_1 \to K_1 \to 0, \ldots$, we have $L_0 \otimes k \cong M \otimes k$,
$L_1 \otimes k \cong K_1 \otimes k, \ldots$. Any two minimal resolutions of $M$ are isomorphic as com-
plexes (prove this!).

\begin{example*}
Let $x_1, \ldots, x_n \in m$ be an $A$-sequence, and let $K_\bullet = K_\bullet(x_1, \ldots, x_n)$
be the Koszul complex
\[
  0 \to K_n \to K_{n-1} \to \cdots \to K_0 \to A/(x_1, \ldots, x_n) \to 0;
\]
then $K_\bullet$ is a minimal resolution of $A/(x_1, \ldots, x_n)$ over $A$.
\end{example*}

Let $(A, m, k)$ be a Noetherian local ring; then a finite $A$-module $M$
always has a minimal resolution. Construction: let $\{\omega_1, \ldots, \omega_p\}$ be a
minimal basis of $M$, let $L_0 = Ae_1 + \cdots + Ae_p$ be a free module, and define
$\varepsilon\colon L_0 \to M$ by $\varepsilon(e_i) = \omega_i$; taking $K_1$ to be the kernel of $\varepsilon$ we get
$0 \to K_1 \to L_0 \to M \to 0$ with $L_0 \otimes k \cong M \otimes k$. Now $K_1$ is again a
finite $A$-module, so that we need only proceed as before.

\begin{lemma}\label{lem:19.1}
Let $(A, m, k)$ be a local ring, and $M$ a finite $A$-module. Suppose
that $L_\bullet$ is a minimal resolution of $M$; then
\begin{enumerate}[(i)]
  \item $\dim_k \Tor^A_i(M, k) = \rank L_i$ for all $i$,
  \item $\projdim M = \sup \{i \mid \Tor^A_i(M, k) \neq 0\} \le \projdim_A k$,
  \item if $M \neq 0$ and $\projdim M = r < \infty$ then for any finite $A$-module $N \neq 0$
        we have $\Ext^r_A(M, N) \neq 0$.
\end{enumerate}
\end{lemma}

\begin{proof}
(i) We have $\Tor^A_i(M, k) = H_i(L_\bullet \otimes k)$, but from the definition of
minimal resolution, $\bar{d}_i = 0$, and hence $H_i(L_\bullet \otimes k) = L_i \otimes k$, and the dimen-
sion of this as a $k$-vector space is equal to $\rank_A L_i$.

(ii) follows from (i).

(iii) Since $L_{r+1} = 0$ and $L_r \neq 0$, $\Ext^r_A(M, N)$ is the cokernel of $d_r^*\colon
\Hom_A(L_r, N) \to \Hom_A(L_{r-1}, N)$, but since $L_i$ is free, $\Hom_A(L_i, N)$ is just
a direct sum of a number of copies of $N$; we can write $d_r\colon L_r \to L_{r-1}$ as a
matrix with entries in $m$, and then $d_r^*$ is given by the same matrix, so that
$\Im(d_r^*) \subset m\,\Hom_A(L_r, N)$, and by NAK $\Ext^r_A(M, N) \neq 0$.
\end{proof}

\begin{remark}
One sees from the above lemma that $\Tor_i(M, k) = 0$ implies that
$L_i = 0$, and therefore $\projdim M < i$, so that $\Tor_j(M, k) = 0$ for $j > i$. It is
conjectured that this holds in more generality, or more precisely:
\emph{Rigidity conjecture.}\index{rigidity conjecture, section 19} Let $R$ be a Noetherian ring, $M$ and $N$ finite $R$-modules;
suppose that $\projdim M < \infty$. Then $\Tor^R_i(M, N) = 0$ implies that
$\Tor^R_j(M, N) = 0$ for all $j > i$.
\end{remark}

This has been proved by Lichtenbaum~\cite{Li1} if $R$ is a regular ring, but
is unsolved in general.

The following theorem is not an application of Lemma~\hyperref[lem:19.1]{1}, but is proved
by a similar technique.

\begin{theorem}[Auslander and Buchsbaum]\label{thm:19.1}
Let $A$ be a Noetherian local
ring and $M \neq 0$ a finite $A$-module. Suppose that $\projdim M < \infty$; then
\[
  \projdim M + \depth M = \depth A.
\]
\end{theorem}

\begin{proof}
Set $\projdim M = h$; we work by induction on $h$. If $h = 0$ then $M$
is a free $A$-module, so that the assertion is trivial. If $h = 1$, let
\[
  (\varphi) \quad 0 \to A^m \xrightarrow{\varphi} A^n \to M \to 0
\]
be a minimal resolution of $M$. We can write $\varphi$ as an $m \times n$ matrix with
entries in $m$. From $(\varphi)$ we obtain the long exact sequence
\[
  \cdots \to \Ext^i_A(k, A^n) \xrightarrow{\varphi^*} \Ext^i_A(k, A^m) \to \Ext^{i+1}_A(k, M) \to \cdots,
\]
and writing out $\Ext^i_A(k, A^m) = \Ext^i_A(k, A)^m$ and $\Ext^i_A(k, A^n) = \Ext^i_A(k, A)^n$,
we can express $\varphi^*$ by the same matrix as $\varphi$. However, the entries of $\varphi$ are
elements of $m$, and therefore annihilate $\Ext^i_A(k, A)$, so that $\varphi^* = 0$, and we
have an exact sequence
\[
  0 \to \Ext^i_A(k, A)^n \to \Ext^{i+1}_A(k, M) \to \Ext^{i+1}_A(k, A)^m \to 0
\]
for every $i$. Since $\depth M = \inf \{i \mid \Ext^i_A(k, M) \neq 0\}$ we have $\depth M =
\depth A - 1$ and the theorem holds if $h = 1$. If $h > 1$ then taking any exact
sequence
\[
  0 \to M' \to A^n \to M \to 0,
\]
we have $\projdim M' = h - 1$, so that an easy induction completes the
proof.
\end{proof}

\begin{lemma}\label{lem:19.2}
Let $A$ be a ring and $n > 0$ a given integer. Then the following
conditions are equivalent.
\begin{enumerate}[(1)]
  \item $\projdim M \le n$ for every $A$-module $M$;
  \item $\projdim M \le n$ for every finite $A$-module $M$;
  \item $\injdim N \le n$ for every $A$-module $N$;
  \item $\Ext^{n+1}_A(M, N) = 0$ for all $A$-modules $M$ and $N$.
\end{enumerate}
\end{lemma}

\begin{proof}
$(1) \implies (2)$ is trivial.

$(2) \implies (3)$ For any ideal $I$, the $A$-module $A/I$ is finite, so that
$\Ext^{n+1}_A(A/I, N) = 0$, and hence by \S\ref{sec:18}, Lemma~\hyperref[lem:18.1]{1}, $\injdim N \le n$.

$(3) \implies (4)$ is trivial, and $(4) \implies (1)$ is well-known (see p.~280).
\end{proof}

We define the \emph{global dimension}\index{global dimension} of a ring $A$ by
\[
  \gldim A = \sup \{\projdim M \mid M \text{ is an $A$-module}\}.
\]
According to Lemma~\hyperref[lem:19.2]{2} above, this is also equal to the maximum projective
dimension of all finite $A$-modules. If $(A, m, k)$ is a Noetherian local ring then by
Lemma~\hyperref[lem:19.1]{1}, $\gldim A = \projdim_A k$.

We have defined regular local rings (see \S\ref{sec:14}) as Noetherian local rings
for which $\dim A = \embdim A$, and we have seen that they are integral
domains (Theorem~\ref{thm:14.3}) and CM rings (Theorem~\ref{thm:17.8}). A regular local
ring is Gorenstein (Theorem~\ref{thm:18.1}, (5')). A necessary and sufficient condition
for a Noetherian local ring $(A, m, k)$ to be regular is that $\gr_m(A)$ is a poly-
nomial ring over $k$ (Theorem~\ref{thm:17.10}). The following theorem gives another
important necessary and sufficient condition.

\begin{theorem}[Serre]\label{thm:19.2}
Let $A$ be a Noetherian local ring. Then
\[
  A \text{ is regular} \iff \gldim A = \dim A \iff \gldim A < \infty.
\]
\end{theorem}

\begin{proof}
(I) Suppose that $(A, m, k)$ is an $n$-dimensional regular local ring.
Let $x_1, \ldots, x_n$ be a regular system of parameters; then since this is an
$A$-sequence, the Koszul complex $K_\bullet(x_1, \ldots, x_n)$ is a minimal free resolution of
$A/(x_1, \ldots, x_n) = k$, and $K_n \neq 0$, $K_{n+1} = 0$, so that as we have already seen,
$\gldim A = \projdim k = n$.

(II) Conversely, suppose that $\gldim A = r < \infty$ and $\embdim A = s$. We
prove that $A$ is regular by induction on $s$; we can assume that $s > 0$, that
is $m \neq 0$. Then $m \notin \Ass(A)$: for if $0 \neq a \in A$ is such that $ma = 0$, consider
a minimal resolution
\[
  0 \to L_r \to L_{r-1} \to \cdots \to L_0 \to k \to 0
\]
of $k$ (with $r > 0$); then $L_r \subset mL_{r-1}$, but then $aL_r = 0$, which contradicts the
assumption that $L_r$ is a free module. Thus we can choose $x \in m$ not
contained in $m^2$ or in any associated prime of $A$. Then $x$ is $A$-regular,
hence also $m$-regular, so that if we set $B = A/xA$ then according to Lemma~\hyperref[lem:18.2]{2}
of \S\ref{sec:18}, $\Ext^q_A(m, N) = \Ext^q_B(m/xm, N)$ for all $B$-modules $N$, and hence
we obtain $\projdim_B m/xm \le r$.

Now we prove that the natural map $m/xm \to m/xA$ splits, so that
$m/xA$ is isomorphic to a direct summand of $m/xm$. Since $x \notin m^2$, we can take
a minimal basis $x_1 = x, x_2, \ldots, x_s$ of $m$ starting with $x$ (here $s = \embdim
A$). We set $b = (x_2, \ldots, x_s)$, so that by the minimal basis condition,
$b \cap xA \subset xm$, and therefore there exists a chain
\[
  m/xA = (b + xA)/xA \cong b/(b \cap xA) \to m/xm \to m/xA
\]
of natural maps, whose composite is the identity. This proves the above
claim. Now clearly,
\[
  \projdim_B m/xA \le \projdim_B m/xm \le r.
\]
Taking a minimal $B$-projective resolution of $m/xA$ and patching it together
with the exact sequence $0 \to m/xA \to B \to k \to 0$ gives a projective
resolution of $k$ of length $\le r + 1$, and hence $\gldim B = \projdim_B k \le r + 1$,
so that by induction, $B$ is a regular local ring. Since $x$ is not contained
in any associated prime of $A$ we have $\dim B = \dim A - 1$, and therefore
$A$ is regular.
\end{proof}

\begin{theorem}[Serre]\label{thm:19.3}
Let $A$ be a regular local ring and $P$ a prime ideal; then
$A_P$ is a regular local ring.
\end{theorem}

\begin{proof}
Since $\projdim_A A/P \le \gldim A < \infty$, as an $A$-module $A/P$ has a
projective resolution $L_\bullet$ of finite length. Then $L_\bullet \otimes_A A_P$ is a projective
resolution of $(A/P) \otimes_A A_P = A_P/PA_P = K(P)$ as an $A_P$-module, so that $K(P)$
has a projective resolution of finite length as an $A_P$-module, which means that $A_P$
has finite global dimension; thus by the previous theorem, $A_P$ is
regular.
\end{proof}

\begin{definition*}
A \emph{regular ring}\index{regular rings} is a Noetherian ring such that the localisation
at every prime is a regular local ring. By the previous theorem, it is
enough\index{regular local rings} for the localisation at every maximal ideal to be regular.
\end{definition*}

\begin{theorem}\label{thm:19.4}
A regular ring is normal.
\end{theorem}

\begin{proof}
The definition of normal is local, so that it is enough to show that a
regular local ring is normal. We show that the conditions of the corollary
to Theorem~\ref{thm:11.5} are satisfied. (a) The localisation at a height 1 prime ideal is a
DVR by the previous theorem and Theorem~\ref{thm:11.2}. (b) All the prime
divisors of a non-zero principal ideal have height 1 by Theorem~\ref{thm:17.8} (the
implication regular $\Rightarrow$ CM).
\end{proof}

\begin{theorem}\label{thm:19.5}
If $A$ is regular then so are $A[X]$ and $A[[X]]$.
\end{theorem}

\begin{proof}
For $A[X]$, let $P$ be a maximal ideal of $A[X]$ and set $P \cap A = m$.
Then $A[X]_P$ is a localisation of $A_m[X]$, so that replacing $A$ by $A_m$, we can assume
$A$ is a regular local ring. Then setting $A/m = k$ we have $A[X]/m[X] =
k[X]$, so that there is a monic polynomial $f(X)$ with coefficients in $A$ such that
$P = (m, f(X))$, and such that $f$ reduces to an irreducible polynomial in
$k[X]$ modulo $m$. Then by Theorem~\ref{thm:15.1}, we clearly have
\[
  \dim A[X]_P = \htop P = 1 + \htop m = 1 + \dim A;
\]
on the other hand $m$ is generated by $\dim A$ elements, so that $P = (m, f)$ is
generated by $\dim A + 1$ elements, and therefore $A[X]_P$ is regular.

For $A[[X]]$, set $B = A[[X]]$ and let $M$ be a maximal ideal of $B$; then $X \in M$
by Theorem~\ref{thm:8.2}, (i). Therefore $M \cap A = m$ is a maximal ideal of $A$. Now
although we cannot say that $B_M$ contains $A_m[[X]]$, the two have the same
completion, $(B_M)^\wedge = (A_m)^\wedge[[X]]$. A Noetherian local ring is regular if and
only if its completion is regular (since both the dimension and embedding
dimension remain the same on taking the completion). Thus if we replace
$A$ by $(A_m)^\wedge$, the maximal ideal of $B = A[[X]]$ is $M = (m, X)$, and $\htop M =
\htop m + 1$, so that $B$ is also regular.
\end{proof}

Next we discuss the properties of modules which have finite free
resolutions (the definition is given below).

\begin{lemma}[Schanuel]\label{lem:19.3}
Let $A$ be a ring and $M$ an $A$-module. Suppose that
\[
  0 \to K \to P \to M \to 0 \quad \text{and} \quad 0 \to K' \to P' \to M \to 0
\]
are exact sequences with $P$ and $P'$ projective. Then $K \oplus P' \cong K' \oplus P$.
\end{lemma}

% 注: 原书此处用 λ, λ' 与矩阵 (φ, ψ, φψ = 1) 的写法, 扫描版矩阵不可辨;
% 此处按标准 fibre product 证法转录, 结论与原书一致.
\begin{proof}
From the fact that $P$ and $P'$ are projective, there exist $\lambda\colon P \to P'$
and $\lambda'\colon P' \to P$, giving the diagram
\[
  \begin{array}{ccccccc}
    0 & \to & K & \to & P & \xrightarrow{\varepsilon} & M \to 0 \\
    & & & & \downarrow \lambda' & & \downarrow \lambda \\
    0 & \to & K' & \to & P' & \xrightarrow{\varepsilon'} & M \to 0
  \end{array}
\]
with $\varepsilon'\lambda = \varepsilon$ and $\varepsilon\lambda' = \varepsilon'$.
We add in harmless summands $P'$ and $P$ to the two exact rows, and line
up the middle terms:
\[
  0 \to K \oplus P' \to P \oplus P' \xrightarrow{(\varepsilon, 0)} M \to 0
  \quad \text{and} \quad
  0 \to P \oplus K' \to P \oplus P' \xrightarrow{(0, \varepsilon')} M \to 0.
\]
Consider the submodule $W = \{(x, y) \in P \oplus P' \mid \varepsilon(x) = \varepsilon'(y)\}$ of $P \oplus P'$,
the fibre product of $P$ and $P'$ over $M$. Projection onto the first factor gives
the exact sequence $0 \to K' \to W \to P \to 0$: indeed $W \to P$ is surjective since
$\varepsilon'$ is surjective, and its kernel is $\{(0, y) \mid \varepsilon'(y) = 0\} = K'$. Since $P$ is
projective, this sequence splits, so that $W \cong K' \oplus P$. Similarly, projection onto
the second factor gives the exact sequence $0 \to K \to W \to P' \to 0$, which splits, so
that $W \cong K \oplus P'$. Therefore $\varphi$ induces an isomorphism
$K \oplus P' \cong P \oplus K'$.
\end{proof}

\begin{lemma}[generalised Schanuel lemma]\label{lem:19.4}
Let $A$, $M$ be as above, and suppose
that
\[
  0 \to P_n \to \cdots \to P_1 \to P_0 \to M \to 0
  \quad \text{and} \quad
  0 \to Q_n \to \cdots \to Q_1 \to Q_0 \to M \to 0
\]
are exact sequences with $P_i$ and $Q_i$ projective for
$0 \le i \le n - 1$. Then
\[
  P_0 \oplus Q_1 \oplus P_2 \oplus \cdots \cong Q_0 \oplus P_1 \oplus Q_2 \oplus \cdots
\]
(the alternating sums of $P_i$ and $Q_i$).
\end{lemma}

\begin{proof}
Write $K$ for the kernel of $P_0 \to M$ and $K'$ for the kernel of
$Q_0 \to M$; then, by the previous lemma, $K \oplus Q_0 \cong P_0 \oplus K'$. Now add
in harmless summands $Q_n$ and $P_n$ to $0 \to P_n \to \cdots \to P_1 \to
K \to 0$ and $0 \to Q_n \to \cdots \to Q_1 \to K' \to 0$ respectively, to obtain
exact sequences
\[
  0 \to P_n \to \cdots \to P_1 \oplus Q_0 \to K \oplus Q_0 \to 0
  \quad \text{and} \quad
  0 \to Q_n \to \cdots \to Q_1 \oplus P_0 \to P_0 \oplus K' \to 0,
\]
and induction on $n$ gives the assertion.
\end{proof}

A \emph{finite free resolution} (or \emph{FFR}\index{FFR} for short) of an $A$-module $M$ is a
resolution
\[
  0 \to F_r \to \cdots \to F_1 \to F_0 \to M \to 0 \quad \text{(of finite length)}
\]
with each $F_i$ a finite free module. If $M$ has an FFR we set
\[
  \chi(M) = \sum_{i=0}^{r} (-1)^i \rank F_i,
\]
and call $\chi(M)$ the \emph{Euler number}\index{Euler number} of $M$. By Lemma~\hyperref[lem:19.4]{4},
$\chi(M)$ is independent of the choice of FFR. Moreover, since for any prime
ideal $P$,
\[
  0 \to (F_r)_P \to \cdots \to (F_1)_P \to (F_0)_P \to M_P \to 0
\]
is an FFR of the $A_P$-module $M_P$ we have $\chi(M) = \chi(M_P)$. If $M$ is itself
free then one sees easily from Lemma~\hyperref[lem:19.4]{4} that $\chi(M) = \rank M$.

\begin{theorem}\label{thm:19.6}
Let $(A, m)$ be a local ring, and suppose that for any finite
subset $E \subset m$ there exists $0 \neq y \in A$ such that $yE = 0$; then the only
$A$-modules having an FFR are the free modules.
\end{theorem}

\begin{remark}
If $A$ is Noetherian then the assumption on $m$ is equivalent to
$m \in \Ass(A)$, or $\depth A = 0$. In this case the theorem is a special case of
Theorem~\ref{thm:19.1}.
\end{remark}

\begin{proof}
Suppose that $0 \to F_r \to F_{r-1} \to \cdots \to F_0 \to M \to 0$ is an FFR
of $M$, and set $N = \Coker(F_r \to F_{r-1})$; if we prove that $N$ is free then
we can decrease $r$ by 1, so that we only need consider the case $0 \to
F_1 \to F_0 \to M \to 0$. Now let $0 \to L_1 \to L_0 \to M \to 0$ be a minimal
free resolution of $M$; then since $L_0$ and $F_0$ are finitely generated, by
Schanuel's lemma (or by Theorem~\ref{thm:2.6}), $L_1$ is also finite. Considering bases of
$F_1$ and $L_1$, we can write down a set of generators of $L_1$ as a submodule of
$F_0$ using only a finite number of elements of $m$. Then by assumption, there
exists $0 \neq y \in A$ such that $yL_1 = 0$. Since $L_1$ is a free module, we must
have $L_1 = 0$, so that $M \cong L_0$, and is free.
\end{proof}

\begin{theorem}\label{thm:19.7}
Let $A$ be any ring; if $M$ is an $A$-module having an FFR then
$\chi(M) \ge 0$.
\end{theorem}

\begin{proof}
Choose a minimal prime ideal $P$ of $A$; since $\chi(M) = \chi(M_P)$, we can
replace $A$ by $A_P$, and then $A$ is a local ring with maximal ideal $m$ equal to
$\nil(A)$. Then the assumption of the previous theorem is satisfied: for
$x_1, \ldots, x_r \in m$, we can assume by induction on $r$ that there is a $z \neq 0$
such that $zx_1 = \cdots = zx_{r-1}$; but $x_r$ is nilpotent, so that there is an $i \ge 0$
such that $zx_r^i \neq 0$ but $zx_r^{i+1} = 0$, and we can take $y = zx_r^i$. Therefore by
the previous theorem $M$ is a free module, and $\chi(M) = \rank M \ge 0$.
\end{proof}

\begin{theorem}[Auslander and Buchsbaum~\cite{AB2}]\label{thm:19.8}
Let $A$ be a Noetherian
ring and $M$ an $A$-module, and suppose that $M$ has an FFR. Then the
following three conditions are equivalent:
\begin{enumerate}[(1)]
  \item $\ann(M) \neq 0$;
  \item $\chi(M) = 0$;
  \item $\ann(M)$ contains an $A$-regular element.
\end{enumerate}
\end{theorem}

\begin{proof}
$(1) \implies (2)$ Suppose that $\chi(M) > 0$; then for any $P \in \Ass(A)$ we have
$\chi(M_P) > 0$, and hence $M_P \neq 0$. By Theorem~\hyperref[thm:19.6]{6}, $M_P$ is a free $A_P$-module, so
that setting $I = \ann(M)$ we have $I_P = \ann(M_P) = 0$. If we set $J = \ann(I)$
then this is equivalent to $J \not\subset P$. Since this holds for every $P \in \Ass(A)$
we see that $J$ contains an $A$-regular element, but then $J \cdot I = 0$ implies that
$I = 0$, a contradiction.

$(2) \implies (3)$ If $\chi(M) = 0$ then by Theorem~\hyperref[thm:19.6]{6}, $M_P = 0$ for every $P \in \Ass(A)$.
This means that $\ann(M) \not\subset P$, so that $\ann(M)$ contains an $A$-regular
element.

$(3) \implies (1)$ is obvious.
\end{proof}

\begin{theorem}[Vasconcelos~\cite{Vas1}]\label{thm:19.9}
Let $A$ be a Noetherian local ring, and
$I$ a proper ideal of $A$; assume that $\projdim I < \infty$. Then
\[
  I \text{ is generated by an $A$-sequence} \iff I/I^2 \text{ is a free module over } A/I.
\]
\end{theorem}

\begin{proof}
$(\implies)$ is already known (Theorem~\ref{thm:16.2}). In fact, $I^v/I^{v+1}$ is a free
$A/I$-module for $v = 1, 2, \ldots$

$(\Longleftarrow)$ We can assume that $I \neq 0$. Since $I$ has finite projective dimension
over $A$ so has $A/I$, and since $A$ is local, $A/I$ has an FFR. Now $\ann(A/I) = I$,
so that by the previous theorem $I$ is not contained in any associated prime
of $A$, and therefore we can choose an element $x \in I$ such that $x$ is not
contained in $m^2$ or in any associated prime of $A$. Then $x$ is $A$-regular,
and $\bar{x} = x \bmod I^2$ is a member of a basis of $I/I^2$ over $A/I$; let $\bar{x}$,
$\bar{y}_2, \ldots, \bar{y}_n \in I$ be such that their images form a basis of $I/I^2$. Then if we
set $B = A/xA$, we see by the same argument as in (II) of the proof of
Theorem~\hyperref[thm:19.2]{2} that $\projdim_B I/xI < \infty$, and that $I/xA$ is isomorphic to a direct
summand of $I/xI$. We now set $I^* = I/xA$, so that $\projdim_A I^* < \infty$. But
on the other hand one sees easily that $I^*/I^{*2}$ is a free module over $B/I^*$,
and an induction on the number of generators of $I$ completes the proof.
\end{proof}

\begin{remark}
In Lech~\cite{Le1}, a set $x_1, \ldots, x_n$ of elements of $A$ is defined to be
\emph{independent}\index{independent elements} if
\[
  \sum a_ix_i = 0 \Rightarrow a_i \in (x_1, \ldots, x_n) \quad \text{for all } i.
\]
If we set $I = (x_1, \ldots, x_n)$ then this condition is equivalent to saying that
the images of $x_1, \ldots, x_n$ in $I/I^2$ form a basis of $I/I^2$ over $A/I$. Then if $A$
and $I$ satisfy the hypotheses of the previous theorem, the theorem tells us
that $I = (y_1, \ldots, y_n)$ with $y_1, \ldots, y_n$ an $A$-sequence. Setting $x_i = \sum a_{ij}y_j$ we
see that the matrix $(a_{ij})$ is invertible when considered in $A/I$; this means that
the determinant of $(a_{ij})$ is not in the maximal ideal of $A$, and so $(a_{ij})$ itself
is invertible. Thus $x_1, \ldots, x_n$ is an $A$-quasi-regular sequence, hence an
$A$-sequence. In particular, we get the following corollary.
\end{remark}

\begin{corollary*}
Let $(A, m)$ be a regular local ring. Then if $x_1, \ldots, x_n \in m$ are
independent in the sense of Lech, they form an $A$-sequence.
\end{corollary*}

However, if we try to prove this corollary as it stands, the induction
does not go through. The key to success with Vasconcelos' theorem is to
strengthen the statement so that induction can be used effectively. Now
as Kaplansky has also pointed out, the main part of Theorem~\hyperref[thm:19.2]{2} (the
implication $\gldim A < \infty \Rightarrow$ regular) follows at once from Theorem~\hyperref[thm:19.9]{9},
because if $m$ is generated by an $A$-sequence then $\embdim A \le \depth A \le
\dim A$.

% ============ 习题 ============
\section*{Exercises to \S\ref{sec:19}}
\emph{Prove the following propositions.}

\begin{exercise}\label{ex:19.1-ex}
Let $k$ be a field and $R = R_0 + R_1 + R_2 + \cdots$ a Noetherian graded ring
with $R_0 = k$; set $m = R_1 + R_2 + \cdots$. Show that if $R_m$ is an $n$-dimensional
regular local ring then $R$ is a polynomial ring $R = k[y_1, \ldots, y_n]$ with $y_i$
homogeneous of positive degree.
\end{exercise}

\begin{exercise}\label{ex:19.2-ex}
Let $A$ be a ring and $M$ an $A$-module. Say that $M$ is \emph{stably free}\index{stably free} if there exist
finite free modules $F$ and $F'$ such that $M \oplus F = F'$. Obviously a stably free
$A$-module $M$ is a finite projective $A$-module, and has an FFR $0 \to
F \to F' \to M \to 0$. Prove that, conversely, a finite projective module
having an FFR is stably free.
\end{exercise}

\begin{exercise}\label{ex:19.3-ex}
Prove that if every finite projective module over a Noetherian ring $A$ is
stably free then every finite $A$-module of finite projective dimension has an
FFR.
\end{exercise}

\begin{exercise}\label{ex:19.4-ex}
Prove that if every finite module over a Noetherian ring $A$ has an FFR
then $A$ is regular.
\end{exercise}

% ============ §20 UFDs (原书 pp.161-169) ============
\section{UFDs}\label{sec:20}\index{UFD, section 20}\index{factorial ring, section 20}

This section treats UFDs, which we have already touched on in
\S\ref{sec:1}; note that the Bourbaki terminology for UFD is `factorial ring'. First of
all, we have the following criterion for Noetherian rings.

\begin{theorem}\label{thm:20.1}
A Noetherian integral domain $A$ is a UFD if and only if
every height 1 prime ideal is principal.
\end{theorem}

\begin{proof}
Proof of `only if'. Suppose that $A$ is a UFD and that $P$ is a height 1
prime ideal. Take any non-zero $a \in P$, and express $a$ as a product of prime
elements, $a = \prod \pi_i$. Then at least one of the $\pi_i$ belongs to $P$; if $\pi_i \in P$ then
$(\pi_i) \subset P$, but $(\pi_i)$ is a non-zero prime ideal and $\htop P = 1$, hence $P = (\pi_i)$.

Proof of `if'. Since $A$ is Noetherian, every element $a \in A$ which is neither
0 nor a unit can be written as a product of finitely many irreducibles.
Hence it will be enough to prove that an irreducible element $a$ is a prime
element. Let $P$ be a minimal prime divisor of $(a)$; then by the principal
ideal theorem (Theorem~\ref{thm:13.9}), $\htop P = 1$, so that by assumption we can
write $P = (b)$. Thus $a = bc$, and since $a$ is irreducible, $c$ is a unit, so that
$(a) = (b) = P$, and $a$ is a prime element.
\end{proof}

\begin{theorem}\label{thm:20.2}
Let $A$ be a Noetherian integral domain, $\Gamma$ a set of prime
elements of $A$, and let $S$ be the multiplicative set generated by $\Gamma$. If $A_S$
is a UFD then so is $A$.
\end{theorem}

\begin{proof}
Let $P$ be a height 1 prime ideal of $A$. If $P \cap S \neq 0$ then $P$ contains
an element $\pi \in \Gamma$, and since $\pi A$ is a non-zero prime ideal we have $P = \pi A$.
If $P \cap S = 0$ then $PA_S$ is a height 1 prime ideal of $A_S$, so that $PA_S = aA_S$
for some $a \in P$. Among all such $a$ choose one such that $aA$ is maximal;
then $a$ is not divisible by any $\pi \in \Gamma$. Now if $x \in P$ we have $sx = ay$ for
some $s \in S$ and $y \in A$. Let $s = \pi_1 \cdots \pi_r$ with $\pi_i \in \Gamma$; then $a \notin \pi_iA$, so
that $y \in \pi_iA$, and an induction on $r$ shows that $y \in sA$, so that $x \in aA$. Hence
$P = aA$.
\end{proof}

\begin{lemma}\label{lem:20.1}
Let $A$ be an integral domain, and $a$ an ideal of $A$ such that
$a \oplus A^n \cong A^{n+1}$; then $a$ is principal.
\end{lemma}

\begin{proof}
Fix the basis $e_0, \ldots, e_n$ of $A^{n+1}$, and viewing $a \oplus A^n \subset A \oplus A^n$, fix
$f_0, \ldots, f_n$ such that $f_0$ is a basis of $A$ and $f_1, \ldots, f_n$ a basis of $A^n$. Then
the isomorphism $\varphi\colon A^{n+1} \to a \oplus A^n$ can be given in the form $\varphi(e_i) =
\sum_{j=0}^{n} a_{ij}f_j$. Write $d_i$ for the $(i, 0)$th cofactor of the matrix $(a_{ij})$, and $d$
for the determinant, so that, since $\varphi$ is injective, $d \neq 0$, and $\sum_i a_{i0}d_i = d$,
$\sum_i a_{ij}d_i = 0$ if $j \neq 0$. Hence if we set $e_0' = \sum_i d_ie_i$, we have $\varphi(e_0') = df_0$.
Moreover, since the image of $\varphi$ includes $f_1, \ldots, f_n$, there exist $e_1', \ldots,
e_n' \in A^{n+1}$ such that $\varphi(e_j') = f_j$. Now define a matrix $(c_{jk})$ by $e_j' = \sum_{k=0}^{n} c_{jk}e_k$
for $j = 0, \ldots, n$ (so $c_{0k} = d_k$). Then we have
\[
  (c_{jk})(a_{ij}) = \begin{pmatrix}
    d & 0 & \cdots & 0 \\
    0 & 1 & \cdots & 0 \\
    \vdots & \vdots & \ddots & \vdots \\
    0 & 0 & \cdots & 1
  \end{pmatrix},
\]
so that by comparing the determinants of both sides we get $\det(c_{jk}) = 1$.
Therefore $e_0', \ldots, e_n'$ is another basis of $A^{n+1}$, and $\varphi(Ae_0') = dAf_0$,
so that $a = dA$.
\end{proof}

Let $K$ be the field of fractions of the integral domain $A$; for a finite
$A$-module $M$, the dimension of $M \otimes_A K$ as a vector space over $K$ is called
the \emph{rank}\index{rank of a module} of $M$. A torsion-free finite $A$-module of rank 1 is isomorphic to
an ideal of $A$. Lemma~\hyperref[lem:20.1]{1} can be formulated as saying that for an integral
domain $A$, a stably free rank 1 module is free (see Ex.~\ref{ex:19.2-ex}). The elementary
proof given above is taken from a lecture by M. Narita in 1971.

\begin{theorem}[Auslander and Buchsbaum~\cite{AB3}]\label{thm:20.3}
A regular local ring is a
UFD.
\end{theorem}

\begin{proof}
Let $(A, m)$ be a regular local ring: the proof works by induction
on $\dim A$. If $\dim A = 0$ then $A$ is a field and therefore (trivially) a UFD.
If $\dim A = 1$ then $A$ is a DVR, and therefore a UFD. We suppose that
$\dim A > 1$ and choose $x \in m - m^2$; then since $xA$ is a prime ideal, applying
Theorem~\hyperref[thm:20.2]{2} to $\Gamma = \{x\}$, we need only show that $A_x$ is a UFD (where
$A_x = A[x^{-1}]$ is as on p.~22). Let $P$ be a height 1 prime ideal of $A_x$, and
set $p = P \cap A$; we have $P = pA_x$. Since $A$ is a regular local ring, the
$A$-module $p$ has an FFR, so that the $A_x$-module $P$ has an FFR. For any
prime ideal $Q$ of $A_x$, the ring $(A_x)_Q = A_{Q \cap A}$ is a regular local ring of
dimension less than that of $A$, so by induction is a UFD. Thus $P_Q$ is free
as an $(A_x)_Q$-module, so that by Theorem~\ref{thm:7.12}, the $A_x$-module $P$ is projective;
hence by Ex.~\ref{ex:19.2-ex}, $P$ is stably free, and therefore by the previous lemma,
$P$ is a principal ideal of $A_x$.
\end{proof}

The above proof is due to Kaplansky. Instead of our Lemma~\hyperref[lem:20.1]{1}, he used
the following more general proposition, which he had previously proved:
if $A$ is an integral domain, and $I_i, J_i$ are ideals of $A$ for $1 \le i \le r$ such that
$\bigoplus_{i=1}^{r} I_i \cong \bigoplus_{i=1}^{r} J_i$, then $I_1 \cdots I_r \cong J_1 \cdots J_r$. This is an interesting property
of ideals, and we have given a proof in Appendix C.

\begin{theorem}\label{thm:20.4}
Let $A$ be a Noetherian integral domain. Then if any finite
$A$-module has an FFR, $A$ is a UFD.
\end{theorem}

\begin{proof}
By Ex.~\ref{ex:19.4-ex}, $A$ is a regular ring. Let $P$ be a height 1 prime ideal
of $A$. Then $A_m$ is a regular local ring for any $m \in \Spec A$, so by the previous
theorem, the ideal $P_m$ is principal, and is therefore a free $A_m$-module. Hence
by Theorem~\ref{thm:7.12}, $P$ is projective. Therefore by Ex.~\ref{ex:19.2-ex}, $P$ is stably free,
and so by Lemma~\hyperref[lem:20.1]{1} is principal.
\end{proof}

Let $A$ be an integral domain; for any two non-zero elements $a, b \in A$,
the notion of greatest common divisor (g.c.d.) and least common multiple
(1.c.m.) are defined as in the ring of integers. That is, $d$ is a g.c.d. of $a$ and
$b$ if $d$ divides both $a$ and $b$, and any element $x$ dividing both $a$ and $b$
divides $d$; and $e$ is an 1.c.m. of $a$ and $b$ if $e$ is divisible by both $a$ and $b$,
and any $y$ divisible by $a$ and $b$ is divisible by $e$; this condition is equivalent
to $(e) = (a) \cap (b)$.

\begin{lemma}\label{lem:20.2}
If an 1.c.m. of $a$ and $b$ exists then so does a g.c.d.
\end{lemma}

\begin{proof}
If $(a) \cap (b) = (e)$ then there exists $d$ such that $ab = ed$. From $e \in (a)$
we get $b \in (d)$ and similarly $a \in (d)$, so that $(a, b) \subset (d)$. Now if $x$ is a common
divisor of $a$ and $b$ then $a = xt$ and $b = xs$, so that $xst$ is a common multiple
of $a, b$, and is hence divisible by $e$. Then from $ed = ab = x \cdot xst$ we get that
$d$ is divisible by $x$. Therefore, $d$ is a g.c.d. of $a$ and $b$.
\end{proof}

\begin{remark}\label{rem:20.1}
If $A$ is a Noetherian integral domain which is not a UFD then
$A$ has an irreducible element $a$ which is not prime. If $xy \in (a)$ but $x \notin (a)$,
$y \notin (a)$ then the only common divisors of $a$ and $x$ are units, so that 1 is a
g.c.d. of $a$ and $x$. However, $xy \in (a) \cap (x)$, but $xy \notin (ax)$, so that $(a) \cap
(x) \neq (ax)$, and there does not exist any 1.c.m. of $a$ and $x$. Thus the converse
of Lemma~\hyperref[lem:20.2]{2} does not hold in general.
\end{remark}

\begin{remark}\label{rem:20.2}
If $A$ is a UFD then an intersection of an arbitrary collection
of principal ideals is again principal (we include (0)). Indeed, if $\bigcap_i a_iA \neq 0$,
then factorise each $a_i$ as a product of primes:
\[
  a_i = u_i \prod_p p^{n_{ip}},
\]
with $u_i$ units, and $p, \pi$ prime elements such that $pA \neq \pi A$ for $p \neq \pi$.
(We could even allow the $a_i$ to be elements of the field of fractions of $A$.)
Then $\bigcap_i a_iA = dA$, where $d = \prod_p p^{\min_i n_{ip}}$.
\end{remark}

\begin{theorem}\label{thm:20.5}
An integral domain $A$ is a UFD if and only if the ascending
chain condition holds for principal ideals, and any two elements of $A$ have
an 1.c.m.
\end{theorem}

\begin{proof}
The `only if' is already known, and we prove the `if'. From the first
condition it follows that every element which is neither 0 nor a unit can
be written as a product of a finite number of irreducible elements, so that
we need only prove that an irreducible element is prime. Let $a$ be an
irreducible element, and let $xy \in (a)$ and $x \notin (a)$. By assumption we can
write $(a) \cap (x) = (z)$; now 1 is a g.c.d. of $a$ and $x$, so that one sees from the
proof of Lemma~\hyperref[lem:20.2]{2} that $(z) = (ax)$, and then $xy \in (a) \cap (x) = (ax)$ implies
that $y \in (a)$. Therefore $(a)$ is prime.
\end{proof}

\begin{theorem}\label{thm:20.6}
Let $A$ be a regular ring and $u, v \in A$. Then $uA \cap vA$ is a
projective ideal.
\end{theorem}

\begin{proof}
$A_m$ is a UFD for every maximal ideal $m$, so that $(uA \cap vA)A_m =
uA_m \cap vA_m$ is a principal ideal, and hence a free module.
\end{proof}

\begin{theorem}\label{thm:20.7}
If $A$ is a UFD then a projective ideal is principal.
\end{theorem}

\begin{proof}
By Theorem~\ref{thm:11.3}, it is equivalent to say that a non-zero ideal $a$ is
projective or invertible. Hence if we set $K$ for the field of fractions of $A$,
then there exist $u_i \in K$ such that $u_ia \subset A$ and $a_i \in a$ such that $\sum u_ia_i = 1$.
We have $a \subset \bigcap u_i^{-1}A$, and conversely if $x \in \bigcap u_i^{-1}A$ then $x =
\sum (xu_i)a_i \in a$, and hence $a = \bigcap u_i^{-1}A$; now since $A$ is a UFD, the intersection
of principal fractional ideals is again principal.
\end{proof}

\begin{theorem}\label{thm:20.8}
If $A$ is a regular UFD then so is $A[X]$.
\end{theorem}

\begin{proof}
Set $B = A[X]$. By Theorem~\hyperref[thm:20.5]{5}, it is enough to prove that $uB \cap vB$
is principal for $u, v \in B$; set $a = uB \cap vB$. Then by Theorem~\hyperref[thm:20.6]{6} and
Theorem~\ref{thm:19.5}, $a$ is projective, so that
\[
  a \otimes_A A = a \otimes_A (B/XB) = a/Xa
\]
is projective as an $A$-module. Suppose that $a = X^rb$ with $b \nsubseteq XB$; then
$a/Xa \cong b/Xb$, so that $b$ is isomorphic to $a$, hence projective, and therefore
locally principal. $B$ is a regular ring, so that the prime divisors of $b$ all
have height 1. Since $XB$ is also a height 1 prime ideal and $b \nsubseteq XB$ we
have $b:XB = b$, hence $b \cap XB = Xb$. Therefore since we can view $b/Xb$
as $b/Xb = b/b \cap XB \subset B/XB = A$, by Theorem~\hyperref[thm:20.7]{7} it is principal, hence
$b = yB + Xb$ for some $y \in b$; then by NAK, $b = yB$, so that $a = X^ryB$.
\end{proof}

\begin{remark}
There are examples where $A$ is a UFD but $A[X]$ is not.
\end{remark}

It is easy to see that a UFD is a Krull ring. For any Krull ring $A$ we
can define the \emph{divisor class group}\index{divisor class group} of $A$, which should be thought of as a
measure of the extent to which $A$ fails to be a UFD. We can give the
definition in simple terms as follows: let $\mathscr{P}$ be the set of height 1 prime
ideals of the Krull ring $A$, and $D(A)$ the free Abelian group on $\mathscr{P}$. That
is, $D(A)$ consists of formal sums $\sum_{p \in \mathscr{P}} n_p \cdot p$ (with $n_p \in \mathbb{Z}$ and all but finitely
many $n_p = 0$), with addition defined by
\[
  \left(\sum n_pp\right) + \left(\sum n_p'p\right) = \sum (n_p + n_p')p.
\]
Let $K$ be the field of fractions of $A$, and $K^*$ the multiplicative group of
non-zero elements of $K$, and for $a \in K^*$ set $\divisor(a) = \sum_{p \in \mathscr{P}} v_p(a) \cdot p$, where $v_p$
is the normalised additive valuation of $K$ corresponding to $p$. Then
$\divisor(ab) = \divisor(a) + \divisor(b)$, so that $\divisor$ is a homomorphism from $K^*$ to $D(A)$. We
write $F(A)$ for the image of $K^*$; this is a subgroup of $D(A)$, so that
we can define $C(A) = D(A)/F(A)$ to be the divisor class group of $A$.

Now, if $A$ is a UFD then each $p \in \mathscr{P}$ is principal, and if $p = aA$ then for
an element of $D(A)$ we have $p = \divisor(a)$, so that $C(A) = 0$. Conversely, if
$C(A) = 0$ then each $p \in \mathscr{P}$ is a principal ideal, and putting this together
with the corollary of Theorem~\ref{thm:12.3}, one sees easily that $A$ is a UFD. Hence
\[
  A \text{ is a UFD} \iff C(A) = 0.
\]

Now let $A$ be any ring, and $M$ a finite projective $A$-module. For each
$P \in \Spec A$, the localisation $M_P$ is a free module over $A_P$, and we write
$n(P)$ for its rank. Then $n$ is a function on $\Spec A$, and is constant on every
connected component (since $n(P) = n(Q)$ if $P \supset Q$). This function $n$ is called
the \emph{rank} of $M$. If the rank is a constant $r$ over the whole of $\Spec A$ then
we say that $M$ is a \emph{projective module of rank $r$}\index{projective module of rank r}. We write $\Pit(A)$ for the
set of isomorphism classes of finite projective $A$-modules of rank 1; $\cl(M)$
denotes the isomorphism class of $M$. If $M$ and $N$ are finite projective rank
1 modules then so is $M \otimes_A N$; this is clear on taking localisations. Thus
we can define a sum in $\Pit(A)$ by setting
\[
  \cl(M) + \cl(N) = \cl(M \otimes N).
\]
We set $M^* = \Hom_A(M, A)$, and define $\varphi\colon M \otimes M^* \to A$ by
\[
  \varphi\left(\sum m_i \otimes f_i\right) = \sum f_i(m_i);
\]
then $\varphi$ is an isomorphism (taking localisations and using the corollary to
Theorem~\ref{thm:7.11} reduces to the case $M = A$, which is clear). Hence $\cl(M^*) =
-\cl(M)$, and $\Pit(A)$ becomes an Abelian group, called the \emph{Picard group}\index{Picard group}
of $A$. If $A$ is local then $\Pit(A) = 0$.

If $A$ is an integral domain with field of fractions $K$, then $M_{(0)} = M \otimes K$,
so that the rank we have just defined coincides with the earlier definition
(after Lemma~\hyperref[lem:20.1]{1}). If $M$ is a finite projective rank 1 module, then since $M$
is torsion-free we have $M \subset M_{(0)} \cong K$, so that $M$ is isomorphic as an
$A$-module to a fractional ideal; for fractional ideals, by Theorem~\ref{thm:11.3},
projective and invertible are equivalent conditions, so that for an integral
domain $A$, we can consider $\Pit(A)$ as a quotient of the group of invertible
fractional ideals under multiplication. A fractional ideal $I$ is isomorphic
to $A$ as an $A$-module precisely when $I$ is principal, so that
\[
  \Pit(A) = (\text{invertible fractional ideals})/(\text{principal fractional ideals}).
\]

Suppose in addition that $A$ is a Krull ring. Then we can view $\Pit(A)$
as a subgroup of $C(A)$. To prove this, for $p \in \mathscr{P}$ and $I$ a fractional ideal,
set
\[
  v_p(I) = \min \{v_p(x) \mid x \in I\};
\]
this is zero for all but finitely many $p \in \mathscr{P}$ (check this!), so that we can set
\[
  \divisor(I) = \sum_{p \in \mathscr{P}} v_p(I) \cdot p \in D(A).
\]
For a principal ideal $I = aA$ we have $\divisor(I) = \divisor(a)$. One sees easily that
$\divisor(II') = \divisor(I) + \divisor(I')$, and that $\divisor(A) = 0$, so that if $I$ is invertible,
$\divisor(I) = -\divisor(I^{-1})$.

For invertible $I$ we have $(I^{-1})^{-1} = I$: indeed, $I \subset (I^{-1})^{-1}$ from the
definition, and $I = I \cdot A \supset I(I^{-1}(I^{-1})^{-1}) = (I^{-1})^{-1}$, so $I = (I^{-1})^{-1}$. If $I$ is invertible and
$\divisor(I) = 0$ then $\divisor(I^{-1}) = 0$, so that $I \subset A$, $I^{-1} \subset A$; hence $A \subset (I^{-1})^{-1} = I$,
and $I = A$. It follows that if $I, I'$ are invertible, $\divisor(I) = \divisor(I')$ implies $I = I'$.
Thus we can view the group of invertible fractional ideals as a subgroup of
$D(A)$, and $\Pit(A)$ as a subgroup of $C(A)$.

If $A$ is a regular ring then as we have seen, $p \in \mathscr{P}$ is a locally free module,
and so is invertible. Clearly from the definition, $\divisor(p) = p$. Hence, in
the case of a regular ring, $D(A)$ is identified with the group of invertible
fractional ideals, and $C(A)$ coincides with $\Pit(A)$.

The notions of $D(A)$ and $\Pit(A)$ originally arise in algebraic geometry.
Let $V$ be an algebraic variety, supposed to be irreducible and normal. We
write $\mathscr{W}$ for the set of irreducible codimension 1 subvarieties of $V$, and
define the group of divisors $D(V)$ of $V$ to be the free Abelian group on $\mathscr{W}$;
a divisor (or Weil divisor) is an element of $D(V)$. Corresponding to a
rational function $f$ on $V$ and an element $W \in \mathscr{W}$, let $v_W(f)$ denote the
order of zero of $f$ along $W$, or minus the order of the pole if $f$ has a pole
along $W$. Write $\divisor(f) = \sum_{W \in \mathscr{W}} v_W(f) \cdot W$ for the divisor of $f$ on $V$ (or just
$\divisor(f)$). For $W \in \mathscr{W}$, the local ring $\mathcal{O}_W$ of $W$ on $V$ is a DVR of the function
field of $V$, and $v_W$ is the corresponding valuation. We say that two divisors
$M, N$ in $D(V)$ are linearly equivalent if their difference $M - N$ is the divisor of a
function, and write $M \sim N$. The quotient group of $D(V)$ by $\sim$, that is the
quotient by the subgroup of divisors of functions, is the divisor class group
of $V$ (up to linear equivalence), and we write $C(V)$ for this. (In addition
to linear equivalence one also considers other equivalence relations with
certain geometric significance (algebraic equivalence, numerical equivalence,
$\ldots$), and divisor class groups, quotients of $D(V)$ by the corresponding subgroups.)

A divisor $M$ on $V$ is said to be a \emph{Cartier divisor} if it is the divisor of a
function in a neighbourhood of every point of $V$. From a Cartier divisor $M$
one constructs a line bundle over $V$, and two Cartier divisors give rise to
isomorphic line bundles if and only if they are linearly equivalent. Cartier
divisors form a subgroup of $D(V)$, and their class group up to linear
equivalence is written $\Pit(V)$; this can also be considered as the group of
isomorphism classes of line bundles over $V$ (with group law defined by
tensor product). If $V$ is smooth then (by Theorem~\hyperref[thm:20.3]{3}) there is no distinction
between Cartier and Weil divisors, and $C(V) = \Pit(V)$.

The reader familiar with algebraic geometry will know that the divisor
class group and Picard group of a Krull ring are an exact translation of
the corresponding notions in algebraic geometry. If $V$ is an affine variety,
with coordinate ring $k[V] = A$ then $C(V) = C(A)$ and $\Pit(V) = \Pit(A)$. In
this case, to say that $A$ is a UFD expresses the fact that every codimension
1 subvariety of $V$ can be defined as the intersection of $V$ with a hypersurface.
If $V \subset \mathbb{P}^n$ is a projective algebraic variety, defined by a prime ideal
$I \subset k[X_0, \ldots, X_n]$, and we set $A = k[X]/I = k[t_0, \ldots, t_n]$ (with $t_i$ the
class of $X_i$) then $A$ is the so-called homogeneous coordinate ring of $V$. If
$A$ is integrally closed we say that $V$ is projectively normal (also
arithmetically normal). This condition is stronger than saying that $V$ is
normal (the local ring of any point of $V$ is normal). If $A$ is a UFD then
every codimension 1 subvariety of $V$ can be given as the intersection in
$\mathbb{P}^n$ of $V$ with a hypersurface. Let $m = (t_0, \ldots, t_n)$ be the homogeneous
maximal ideal of $A$, and write $R = A_m$ for the localisation. The above
statement holds if we just assume that $R$ is a UFD; see Ex.~\ref{ex:20.6-ex}. All the
information about $V$ is contained in the local ring $R$.

Thus $C(A)$, $\Pit(A)$ and the UFD condition are notions with important
geometrical meaning, and methods of algebraic geometry can also be used
in their study. For example, in this way Grothendieck~\cite{G5} was able to
prove the following theorem conjectured by Samuel: let $R$ be a regular
local ring, $P$ a prime ideal generated by an $R$-sequence, and set $A = R/P$;
if $A_p$ is a UFD for every $p \in \Spec A$ with $\htop p \le 3$ then $A$ is a UFD.

We do not have the space to discuss $C(A)$ and $\Pit(A)$ in detail, and we
just mention the following two theorems as examples:
\begin{enumerate}[(1)]
  \item If $A$ is a Krull ring then $C(A) \cong C(A[X])$.
        This generalises the well-known theorem (see Ex.~\ref{ex:20.2-ex}) that if $A$ is a UFD
        then so is $A[X]$.
  \item If $A$ is a regular ring then $C(A) \cong C(A[[X]])$.
        This generalises Theorem~\hyperref[thm:20.8]{8}.
\end{enumerate}
Finally we give an example. Let $k$ be a field of characteristic 0, and set
$A = k[X, Y, Z]/(Z^n - XY) = k[x, y, z]$ for some $n > 1$. Then $A/(z, x) \cong
k[X, Y, Z]/(X, Z) \cong k[Y]$, so that $p = (x, z)$ is a height 1 prime ideal of $A$.
In $D(A)$ we have $np = \divisor(x)$, and it can be proved that $C(A) \cong \mathbb{Z}/d$ (see
\cite{S2}, p.~58). The relation $xy = z^n$ shows that $A$ is not a UFD.

For those wishing to know more about UFDs, consult~\cite{K},~\cite{S2} and
\cite{F}.

% ============ 习题 ============
\section*{Exercises to \S\ref{sec:20}}
\emph{Prove the following propositions.}

\begin{exercise}\label{ex:20.1-ex}
(Gauss' lemma) Let $A$ be a UFD, and $f(X) = a_0 + a_1X + \cdots + a_nX^n \in
A[X]$; say that $f$ is \emph{primitive} if the g.c.d. of the coefficients $a_0, \ldots, a_n$
is 1. Then if $f(X)$ and $g(X)$ are primitive, so is $f(X)g(X)$.
\end{exercise}

\begin{exercise}\label{ex:20.2-ex}
If $A$ is a UFD so is $A[X]$ (use the previous question).
\end{exercise}

\begin{exercise}\label{ex:20.3-ex}
If $A$ is a UFD and $q_1, \ldots, q_r$ are height 1 primary ideals then $q_1 \cap \cdots \cap q_r$ is a
principal ideal.
\end{exercise}

\begin{exercise}\label{ex:20.4-ex}
Let $A$ be a Zariski ring (see \S\ref{sec:8}) and $\hat{A}$ the completion of $A$. Then if $A$ is
a UFD so is $\hat{A}$ (there are counter-examples to the converse).
\end{exercise}

\begin{exercise}\label{ex:20.5-ex}
Let $A$ be an integral domain. We say that $A$ is \emph{locally UFD} if $A_m$ is a UFD
for every maximal ideal $m$. If $A$ is a semilocal integral domain and $A$ is
locally UFD, then $A$ is a UFD.
\end{exercise}

\begin{exercise}\label{ex:20.6-ex}
Let $R = \bigoplus_{n \ge 0} R_n$ be a graded ring, and suppose that $R_0$ is a field. Set
$m = \bigoplus_{n > 0} R_n$. If $I$ is a homogeneous ideal of $R$ such that $IR_m$ is principal
then there is a homogeneous element $f \in I$ such that $I = fR$.
\end{exercise}

% ============ §21 Complete intersection rings (原书 pp.169-173) ============
\section{Complete intersection rings}\label{sec:21}\index{complete intersection rings, section 21}\index{c.i. ring, section 21}

Let $(A, m, k)$ be a Noetherian local ring; we choose a minimal basis
$x_1, \ldots, x_n$ of $m$, where $n = \embdim A$ is the embedding dimension (see
\S\ref{sec:14}). Set $E_\bullet = K_\bullet(x_1, \ldots, x_n)$ for the Koszul complex. The complex $E_\bullet$ is
determined by $A$ up to isomorphism. Indeed, if $x_1', \ldots, x_n'$ is another
minimal basis of $m$ then by Theorem~\ref{thm:2.3}, there is an invertible $n \times n$
matrix $(a_{ij})$ over $A$ such that $x_i = \sum a_{ij}x_j'$. It is proved in Appendix C that
$K_\bullet(x_1, \ldots, x_n)$ can be thought of as the exterior algebra $\bigwedge(Ae_1 + \cdots + Ae_n)$ with
differential defined by $d(e_i) = x_i$. Similarly,
$K_\bullet(x_1', \ldots, x_n') = \bigwedge(Ae_1' + \cdots + Ae_n')$ with $d(e_i') = x_i'$.
Now $f(e_j') = \sum a_{ij}e_i$ defines an isomorphism from the free $A$-module
$Ae_1' + \cdots + Ae_n'$ to $Ae_1 + \cdots + Ae_n$, which extends to an isomorphism $f$
of the exterior algebra; $f$ commutes with the differential $d$, since for a
generator $e_i'$ of $\bigwedge(Ae_1' + \cdots + Ae_n')$ we have $df(e_i') = \sum a_{ij}x_j = x_i = fd(e_i')$.
Therefore $f\colon K_\bullet(x_1', \ldots, x_n') \to K_\bullet(x_1, \ldots, x_n)$ is an isomorphism of complexes.

Since $mH_p(E_\bullet) = 0$ by Theorem~\ref{thm:16.4}, $H_p(E_\bullet)$ is a vector space over
$k$, and we set
\[
  \varepsilon_p = \dim_k H_p(E_\bullet) \quad \text{for } p = 0, 1, 2, \ldots;
\]
then these are invariants of the local ring $A$. In view of $H_0(E_\bullet) = A/(x) =
A/m = k$, we have $\varepsilon_0 = 1$. In this section we are concerned with $\varepsilon_1$. If $A$ is
regular then $x_1, \ldots, x_n$ is an $A$-sequence, so that $\varepsilon_1 = \cdots = \varepsilon_n = 0$, and
conversely by Theorem~\ref{thm:16.5}, $\varepsilon_n = 0$ implies that $A$ is regular.

Let us consider the case when $A$ can be expressed as a quotient of a
regular local ring $R$; let $A = R/a$, and write $n$ for the maximal ideal of $R$.
If $a \nsubseteq n^2$, we can take $x \in a - n^2$; then $R' = R/xR$ is again a regular local
ring, and $A = R'/a'$, so that we can write $A$ as a quotient of a ring of
dimension smaller than $R$. In this way we see that there exists an
expression $A = R/a$ of $A$ as a quotient of a regular local ring $(R, n)$ with
$a \subset n^2$. Then we have $m = n/a$ and $m/m^2 = n/(a + n^2) = n/n^2$, so that
$\dim R = \embdim A$. Conversely, equality here implies that $a \subset n^2$.

Let $(R, n)$ be a regular local ring and $A = R/a$ with $a \subset n^2$; choose a
regular system of parameters (that is a minimal basis of $n$) $t_1, \ldots, t_n$.
Then the images $x_i$ of $t_i$ in $A$ form a minimal basis $x_1, \ldots, x_n$ of $m$. Let
\[
  K_\bullet(t_1, \ldots, t_n)\colon 0 \to L_n \to L_{n-1} \to \cdots \to L_1 \to L_0 \to 0
\]
be the Koszul complex of $R$ and $t$. By Theorem~\ref{thm:16.5}, we know that this
becomes exact on adding $\cdots + L_0 \to k \to 0$ to the right-hand end, so that
$K_\bullet(t_1, \ldots, t_n)$ is a projective resolution of $k$ as an $R$-module. Taking the tensor
product with $A = R/a$, we get the complex $E_\bullet = K_\bullet(x_1, \ldots, x_n)$ of $A$-modules.
Thus we have
\[
  H_p(E_\bullet) = H_p(K_\bullet(t_1, \ldots, t_n) \otimes_R A) = \Tor_p^R(k, A) \quad \text{for all } p \ge 0.
\]
However, from the exact sequence of $R$-modules $0 \to a \to R \to A \to 0$ we
get the long exact sequence
\[
  0 = \Tor_1^R(k, R) \to \Tor_1^R(k, A) \to k \otimes_R a \to k \otimes_R R
      \to k \otimes_R A \to 0;
\]
at the right-hand end we have $k \otimes_R R \cong k \otimes_R A = k$, so that
\[
  \Tor_1^R(k, A) \cong k \otimes_R a = a/na.
\]
Quite generally, we write $\mu(M)$ for the minimum number of generators of
an $R$-module $M$. Then we see that
\[
  \mu(a) = \dim_k H_1(E_\bullet) = \varepsilon_1(A).
\]

\begin{theorem}\label{thm:21.1}
Let $(A, m, k)$ be a Noetherian local ring, and $\hat{A}$ its
completion.
\begin{enumerate}[(i)]
  \item $\varepsilon_p(A) = \varepsilon_p(\hat{A})$ for all $p \ge 0$.
  \item $\varepsilon_1(A) \ge \embdim A - \dim A$.
  \item If $R$ is a regular local ring, $a$ an ideal of $R$ and $A \cong R/a$, then
        $\mu(a) = \dim R - \embdim A + \varepsilon_1(A)$.
\end{enumerate}
\end{theorem}

\begin{proof}
(i) is clear from the fact that a minimal basis of $m$ is a minimal
basis of $m\hat{A}$, so that applying $\otimes_A \hat{A}$ to the complex $E_\bullet$ made from $A$
gives that made from $\hat{A}$. Then since $\hat{A}$ is $A$-flat, $H_p(E_\bullet) \otimes \hat{A} = H_p(E_\bullet \otimes \hat{A})$,
and $mH_p(E_\bullet) = 0$ gives $H_p(E_\bullet) \otimes \hat{A} = H_p(E_\bullet)$.

(ii) If $A$ is a quotient of a regular local ring, then as we have seen above,
there exists a regular ring $(R, n)$ such that $A = R/a$ with $a \subset n^2$, so that
$\varepsilon_1(A) = \mu(a) \ge \htop a = \dim R - \dim A = \embdim A - \dim A$, where the
equality for $\htop a$ comes from Theorem~\ref{thm:17.4}, (i). Now $A$ itself is not necessarily
a quotient of a regular local ring, but we will prove later (see \S29) that $\hat{A}$
always is, and we admit this in the section. Having said this, the two sides
of (ii) are unaltered on replacing $A$ by $\hat{A}$, and the inequality holds for $\hat{A}$.

(iii) Set $n = \rad(R)$. If $a \subset n^2$ then, as we have seen above, $\mu(a) = \varepsilon_1(A)$,
and $\dim R = \embdim A$, so that we are done. If $a \nsubseteq n^2$, take $x \in a - n^2$;
we pass to $R/xR$ and $a/xR$, each of $\dim R$ and $\mu(a)$ decreases by 1, so
induction completes the proof.
\end{proof}

\begin{definition*}
A Noetherian local ring $A$ is a \emph{complete intersection ring}
(abbreviated to \emph{c.i. ring}) if $\varepsilon_1(A) = \embdim A - \dim A$.
\end{definition*}

\begin{theorem}\label{thm:21.2}
Let $A$ be a Noetherian local ring.
\begin{enumerate}[(i)]
  \item $A$ is c.i. $\iff$ $\hat{A}$ is c.i.
  \item Let $A$ be a c.i. ring and $R$ a regular local ring such that $A = R/a$;
        then $a$ is generated by an $R$-sequence. Conversely, if $a$ is an ideal generated
        by an $R$-sequence then $R/a$ is a c.i. ring.
  \item A necessary and sufficient condition for $A$ to be a c.i. ring is that
        the completion $\hat{A}$ should be a quotient of a complete regular local ring
        by an ideal generated by an $R$-sequence.
\end{enumerate}
\end{theorem}

\begin{proof}
(i) is obvious.

(ii) By Theorem~\hyperref[thm:21.1]{1}, (iii), $\mu(a) = \dim R - \embdim A + \varepsilon_1(A)$, and by
Theorem~\ref{thm:17.4}, (i), $\htop a = \dim R - \dim A$, so that $A$ is a c.i. ring is equivalent
to $\htop a = \mu(a)$. But by Theorem~\ref{thm:17.4}, (iii), this is equivalent to $a$ being
generated by an $R$-sequence.

(iii) The sufficiency is clear from (i) and (ii). Necessity follows from the
fact that $\hat{A}$ is a quotient of a complete regular local ring (see \S29), together
with (i) and (ii).
\end{proof}

\begin{theorem}\label{thm:21.3}
A c.i. ring is Gorenstein.
\end{theorem}

\begin{proof}
If $A$ is c.i. then so is $\hat{A}$, and if $\hat{A}$ is Gorenstein then so is $A$, so that we
can assume that $A$ is complete. Then we can write $A = R/a$, where $R$ is a
regular local ring and $a$ is an ideal generated by a regular sequence. Since $R$
is Gorenstein, $A$ is also by Ex.~\ref{ex:18.1-ex}.
\end{proof}

Thus we have the following chain of implications for Noetherian local
rings:
\[
  \text{regular} \implies \text{c.i.} \implies \text{Gorenstein} \implies \text{CM}.
\]
Let $A$ be a c.i. ring, and $p$ a prime ideal of $A$. If $A$ is of the form
$R/(x_1, \ldots, x_r)$, where $R$ is regular and $x_1, \ldots, x_r$ is an $R$-sequence, then
since $A_p$ can be written $A_p = R_p/(x_1, \ldots, x_r)$, where $R_p$ is regular and
$x_1, \ldots, x_r$ is an $R_p$-sequence, it follows that $A_p$ is again a c.i. ring. The
question of deciding whether $A_p$ is still a c.i. ring even if $A$ is not a quotient
of a regular local ring remained unsolved for some time, but was answered
affirmatively by Avramov~\cite{Av1}, making use of Andr\'e's homology theory
[1, 2]. % 注: 原书此处引 [1, 2]; André 在文献表中仅一条论文条目, 故保持纯文本
This theory defines homology and cohomology groups
$H_n(A, B, M)$ and $H^n(A, B, M)$ for $n \ge 0$ associated with a ring $A$, an $A$-algebra $B$
and a $B$-module $M$. The definition is complicated, but in any case these are
$B$-modules having various nice functorial properties. If $A$ is a
Noetherian local ring with residue field $k$ then
\[
  A \text{ is regular} \iff H_2(A, k, k) = 0,
\]
and
\[
  A \text{ is c.i.} \iff H_3(A, k, k) = 0;
\]
for $n \ge 3$ the statements $H_3(A, k, k) = 0$ and $H_n(A, k, k) = 0$ are equivalent.
Thus Andr\'e homology is particularly relevant to the study of regular and
c.i. rings.

% ============ 习题 ============
\section*{Exercises to \S\ref{sec:21}}
\emph{Prove the following propositions.}

\begin{exercise}\label{ex:21.1-ex}
Let $R$ be a regular ring, $I$ an ideal of $R$, and let $A = R/I$; then the subset
$\{p \in \Spec A \mid A_p \text{ is c.i.}\}$ is open in $\Spec A$ (use Theorem~\ref{thm:19.9}).
\end{exercise}

\begin{exercise}\label{ex:21.2-ex}
Let $A$ be a Noetherian local ring with $\embdim A = \dim A + 1$; if $A$ is CM
then it is c.i.
\end{exercise}

\begin{exercise}\label{ex:21.3-ex}
Let $k$ be a field, and set $A = k[X, Y, Z]/(X^2 - Y^2, Y^2 - Z^2, XY, YZ, ZX)$;
then $A$ is Gorenstein but not c.i.
\end{exercise}

